% Source PDF: source/普通高考/2011/2011湖北文.pdf
% Page: 1 (the source histogram is vector artwork; no complete embedded bitmap matched it)
% Grid evidence: tmp/repaint_2011_hubei/p1_grid100.jpg and q5_block_grid50.jpg
% Comparison original side: img_comparison/2011/hubei_liberal/q5_fig1.png
% Elements: horizontal sample-data axis with ticks 0,2,4,6,8,10,12; vertical frequency-density axis;
% five adjacent histogram bars on [2,4), [4,6), [6,8), [8,10), [10,12) with heights
% 0.02, 0.05, 0.15, 0.19, and 0.09; dashed guides at 0.02, 0.05, 0.09, 0.15, and 0.19.
% Relations: each bar height is the frequency density for its 2-unit interval; the [10,12) bar
% has density 0.09, giving frequency 0.09*2*200=36 as required by the question.
% solid segments: bar outlines, axis arrows, tick marks, and the two axes.
% dashed segments: horizontal reference guides at 0.02, 0.05, 0.09, 0.15, and 0.19.
% Label strategy: frequency/interval-width label above the y-axis; sample-data label at the x-axis arrow;
% numeric ticks below/left of axes and no labels placed over bars.
\begin{tikzpicture}[baseline]
\begin{axis}[
    width=7.2cm,
    height=5.0cm,
    axis lines=none,
    clip=false,
    xmin=-0.85,
    xmax=13.3,
    ymin=-0.018,
    ymax=0.215,
    ticks=none,
]
    \draw[dashed,line width=0.45pt] (axis cs:0,0.02) -- (axis cs:8,0.02);
    \draw[dashed,line width=0.45pt] (axis cs:0,0.05) -- (axis cs:8,0.05);
    \draw[dashed,line width=0.45pt] (axis cs:0,0.09) -- (axis cs:10,0.09);
    \draw[dashed,line width=0.45pt] (axis cs:0,0.15) -- (axis cs:8,0.15);
    \draw[dashed,line width=0.45pt] (axis cs:0,0.19) -- (axis cs:10,0.19);

    \addplot[
        ybar interval,
        draw=black,
        fill=white,
        line width=0.6pt,
    ] coordinates {
        (2,0.02)
        (4,0.05)
        (6,0.15)
        (8,0.19)
        (10,0.09)
        (12,0)
    };

    \draw[->,line width=0.7pt] (axis cs:0,0) -- (axis cs:13.0,0);
    \draw[->,line width=0.7pt] (axis cs:0,0) -- (axis cs:0,0.21);

    \node[anchor=south west,inner sep=0pt,font=\normalsize] at (axis cs:0.18,0.205) {\(\dfrac{\text{频率}}{\text{组距}}\)};
    \node[anchor=west,inner sep=0pt,font=\normalsize] at (axis cs:12.95,-0.001) {样本数据};

    \draw[line width=0.4pt] (axis cs:0,0) -- (axis cs:0,-0.005);
    \node[anchor=north,inner sep=0pt,font=\normalsize] at (axis cs:0,-0.006) {0};
    \draw[line width=0.4pt] (axis cs:2,0) -- (axis cs:2,-0.005);
    \node[anchor=north,inner sep=0pt,font=\normalsize] at (axis cs:2,-0.006) {2};
    \draw[line width=0.4pt] (axis cs:4,0) -- (axis cs:4,-0.005);
    \node[anchor=north,inner sep=0pt,font=\normalsize] at (axis cs:4,-0.006) {4};
    \draw[line width=0.4pt] (axis cs:6,0) -- (axis cs:6,-0.005);
    \node[anchor=north,inner sep=0pt,font=\normalsize] at (axis cs:6,-0.006) {6};
    \draw[line width=0.4pt] (axis cs:8,0) -- (axis cs:8,-0.005);
    \node[anchor=north,inner sep=0pt,font=\normalsize] at (axis cs:8,-0.006) {8};
    \draw[line width=0.4pt] (axis cs:10,0) -- (axis cs:10,-0.005);
    \node[anchor=north,inner sep=0pt,font=\normalsize] at (axis cs:10,-0.006) {10};
    \draw[line width=0.4pt] (axis cs:12,0) -- (axis cs:12,-0.005);
    \node[anchor=north,inner sep=0pt,font=\normalsize] at (axis cs:12,-0.006) {12};

    \draw[line width=0.4pt] (axis cs:0,0.02) -- (axis cs:-0.12,0.02);
    \frequencyylabel[-4pt]{0}{0.02}{0.02}
    \draw[line width=0.4pt] (axis cs:0,0.05) -- (axis cs:-0.12,0.05);
    \frequencyylabel[-4pt]{0}{0.05}{0.05}
    \draw[line width=0.4pt] (axis cs:0,0.15) -- (axis cs:-0.12,0.15);
    \frequencyylabel[-4pt]{0}{0.15}{0.15}
    \draw[line width=0.4pt] (axis cs:0,0.19) -- (axis cs:-0.12,0.19);
    \frequencyylabel[-4pt]{0}{0.19}{0.19}
\end{axis}
\end{tikzpicture}
